\chapter{Alpha-Fetoprotein (AFP) --- Tumor Marker}

\section*{What Is This Test?}
Alpha-fetoprotein (AFP) is a 70-kDa oncofetal glycoprotein that is the dominant serum protein of the fetus. AFP is synthesized by the fetal yolk sac, fetal liver, and gastrointestinal tract. Its concentration in fetal serum peaks at 12--14 weeks of gestation at levels of 3--4 mg/mL and declines progressively through the remainder of gestation and the first year of life, reaching normal adult levels (<10 ng/mL) by approximately 10--12 months of age. AFP serves as the fetal analog of albumin, functioning as a carrier/transport protein for ligands including fatty acids, bilirubin, copper, and various drugs. In clinical laboratory medicine, AFP is used in three distinct contexts: (1) as a tumor marker for hepatocellular carcinoma (HCC) and non-seminomatous germ cell tumors (NSGCT), (2) as a maternal serum screening test during pregnancy for fetal neural tube defects (open spina bifida, anencephaly) and chromosomal aneuploidy (trisomy 21, trisomy 18), and (3) as a marker for non-malignant liver diseases associated with hepatic regeneration. The AFP tumor marker test discussed here focuses on its role in cancer diagnosis and monitoring.

\section*{Alternative Names}
\begin{itemize}
  \item AFP
  \item Alpha-Fetoprotein
  \item Alpha-Fetoprotein Tumor Marker
  \item AFP Tumor Marker
  \item AFP-L3\% (Lectin-Reactive AFP)
  \item Total AFP
  \item Maternal Serum AFP (MSAFP)
\end{itemize}
\section*{Principle}
AFP is measured by a two-site sandwich chemiluminescent immunoassay (CLIA) or electrochemiluminescent immunoassay (ECLIA) on automated immunoassay analyzers. Two monoclonal antibodies directed against different epitopes on the AFP molecule are used. The capture antibody binds AFP from patient serum; the detection antibody is conjugated to a chemiluminescent label (acridinium ester, ruthenium complex, or alkaline phosphatase). After incubation, washing, and addition of the chemiluminescent substrate, light emission is measured and is proportional to AFP concentration. Results are calibrated against the WHO International Standard for AFP (WHO 1st IS 72/225) and reported in ng/mL. AFP-L3\% (the percentage of total AFP that is reactive to Lens culinaris agglutinin, a fucose-binding lectin) is measured by lectin-affinity electrophoresis followed by immunodetection. AFP synthesized by HCC cells has a higher proportion of fucosylated glycoforms (AFP-L3) compared to AFP from benign hepatocyte regeneration, so an elevated AFP-L3\% ($\geq$10\% of total AFP) increases specificity for HCC over benign liver disease.

\section*{History}
AFP was discovered in 1956 by Bergstrand and Czar in human fetal serum as an alpha-globulin migrating between albumin and alpha-1-globulins on electrophoresis \cite{bergstrand1956demonstration}. In 1963, Yuri Abelev and colleagues demonstrated that AFP was produced by hepatocellular carcinoma in mice, and shortly thereafter, Yury Tatarinov confirmed elevated AFP in human HCC, establishing AFP as the first serum tumor marker \cite{abelev1963production}. In 1972, Brock and Sutcliffe reported elevated AFP in the amniotic fluid of pregnancies affected by fetal neural tube defects \cite{brock1972alpha}. The first AFP immunoassay for clinical use was developed in the 1970s. AFP screening for HCC in high-risk populations (patients with cirrhosis due to hepatitis B or hepatitis C) was widely adopted in the 1990s. The identification of AFP-L3 (the fucosylated isoform of AFP with higher specificity for HCC) by Sato and Taketa in the 1990s led to the development of the AFP-L3\% assay \cite{li2001afpl3}. Current guidelines from the American Association for the Study of Liver Diseases (AASLD) and the European Association for the Study of the Liver (EASL) recommend ultrasound-based screening for HCC every 6 months in at-risk patients, with or without AFP \cite{marrero2018diagnosis}.

\section*{Why This Test Is Ordered}
\begin{itemize}
  \item Screening for hepatocellular carcinoma (HCC) in patients with cirrhosis, particularly from chronic hepatitis B virus (HBV) infection, chronic hepatitis C virus (HCV) infection, alcoholic liver disease, and nonalcoholic steatohepatitis (NASH). AFP measurement every 6 months combined with liver ultrasound is the standard HCC screening protocol, though the value of AFP in addition to ultrasound is debated \cite{marrero2018diagnosis}
  \item Diagnosis and monitoring of non-seminomatous germ cell tumors (NSGCT) of the testis in men and ovarian germ cell tumors in women. AFP is produced by yolk sac (endodermal sinus) tumor elements and embryonal carcinoma. AFP is never elevated in pure seminoma (germinoma) or mature teratoma. Markedly elevated AFP in a testicular mass confirms the presence of a non-seminomatous component and changes staging and treatment
  \item Prognostication and treatment monitoring in HCC --- falling AFP indicates response to locoregional therapy (transarterial chemoembolization [TACE], radiofrequency ablation, Y90-radioembolization) or systemic therapy; rising AFP suggests progression or recurrence. AFP >400 ng/mL in the setting of cirrhosis and a liver mass $>$1 cm has traditionally been considered diagnostic of HCC by imaging alone (LI-RADS criteria)
  \item AFP-L3\% measurement in patients with elevated AFP to distinguish HCC from benign liver disease. An elevated AFP-L3\% ($\geq$10\%) indicates high specificity for HCC \cite{li2001afpl3}
  \item Staging, treatment stratification, and post-treatment monitoring of testicular germ cell tumors along with beta-hCG and LDH, per the International Germ Cell Cancer Collaborative Group (IGCCCG) criteria: AFP >1,000 ng/mL, hCG >5,000 IU/L, or LDH >1.5 times ULN defines intermediate-risk NSGCT; any marker above 10 times ULN defines poor-risk
  \item Monitoring of cirrhosis patients for development of HCC --- serial AFP trends aid in risk stratification, though guidelines emphasize the importance of liver ultrasound
\end{itemize}

\section*{Primary vs Secondary Test}
This is a primary tumor marker test for hepatocellular carcinoma and non-seminomatous germ cell tumors. In the context of HCC screening, AFP is a secondary adjunct to liver ultrasound. In the context of testicular germ cell tumors, AFP is an essential primary tumor marker required for staging, risk classification, and post-treatment monitoring.

\section*{How This Test Is Performed}
For most patients, a health care professional will take a blood sample from a vein in your arm using a small needle. After the needle is inserted, a small amount of blood will be collected into a test tube or vial. You may feel a little sting when the needle goes in or out. This usually takes less than five minutes.

\section*{What Preparation Is Needed}
No special preparation is required. Fasting is not necessary. For testicular germ cell tumor patients, AFP should be measured before orchiectomy whenever possible, as the pre-orchiectomy value is used for IGCCCG risk classification. Serial AFP levels must be measured using the same assay and laboratory for accurate comparison, as AFP immunoassay calibration varies significantly between manufacturers.

\section*{Contraindications and Precautions}
\begin{itemize}
  \item No absolute contraindications. The most important clinical limitation is that AFP is not cancer-specific. AFP is elevated in benign liver diseases associated with hepatocyte necrosis and regeneration: acute hepatitis (any cause --- viral, drug-induced, ischemic), chronic hepatitis (HBV, HCV, autoimmune), and cirrhosis (any cause). In these conditions, AFP typically rises modestly (20--200 ng/mL) and correlates with serum transaminases (ALT/AST), peaking during acute hepatitis flares and declining as transaminases normalize. AFP elevations from benign liver disease are typically transient and correlate with ALT rather than tumor burden. Hepatic metastatic disease from non-HCC primary cancers (particularly gastrointestinal adenocarcinomas metastatic to the liver) can cause AFP elevation from adjacent hepatocyte regeneration rather than AFP production by the tumor. Pregnancy causes physiologic AFP elevation: maternal serum AFP rises from 12 weeks of gestation, peaks at 30--32 weeks, and returns to normal after delivery. Extremely rare causes of elevated AFP: hereditary persistence of AFP (autosomal dominant trait with mildly elevated AFP in healthy individuals without liver disease or malignancy), and ataxia-telangiectasia. Specimen should be collected in a serum separator tube (SST, gold-top) or red-top tube. Approximately 30--40\% of HCC cases (particularly early-stage HCC) have normal AFP (<20 ng/mL).
\end{itemize}
\section*{Risks}
Minimal risk. Slight pain or bruising at the venipuncture site. Rare: hematoma, infection, vasovagal syncope.

\section*{What Happens After the Test}
Results are typically available within 1 to 2 hours. AFP >20 ng/mL in a patient with cirrhosis warrants evaluation for HCC with liver ultrasound or multiphasic CT/MRI. AFP >400 ng/mL in a patient with cirrhosis and a liver mass >1 cm on imaging is highly suspicious for HCC. Markedly elevated AFP (>1,000 ng/mL) in a young man with a testicular mass confirms non-seminomatous germ cell tumor and requires urgent urologic oncology referral. AFP levels should be monitored after curative-intent resection or ablation of HCC or NSGCT: normal post-treatment AFP indicates complete response, while persistent elevation or rising levels indicate residual or recurrent disease.

\section*{Understanding Results}

\subsection*{Below Normal}
\begin{itemize}
  \item AFP <10 ng/mL: Normal level in healthy adults. Does not exclude HCC --- 30--40\% of HCC patients, particularly those with early-stage tumors (<2 cm), have normal AFP levels. Does not exclude NSGCT, as AFP is elevated in only 50--70\% of NSGCT
  \item AFP normalization after curative-intent liver resection, liver transplantation, or locoregional ablation for HCC: Indicates complete response. AFP half-life is 4--6 days; normalization is expected within 4--6 weeks after complete tumor removal
  \item AFP normalization after orchiectomy and/or chemotherapy for NSGCT: Indicates complete clinical response. AFP half-life is 5--7 days. Failure of AFP to normalize according to expected half-life decay indicates residual viable non-seminomatous tumor. If AFP does not decline appropriately after orchiectomy (serum half-life >7 days), there is metastatic disease
  \item Normal AFP in a patient with biopsy-proven HCC that previously had elevated AFP: The tumor may have dedifferentiated and lost AFP production. AFP normalization in this context does not necessarily reflect treatment response. Alternatively, a rising tumor with falling AFP may indicate clonal evolution toward an AFP-non-producing phenotype
  \item Normal AFP in the setting of elevated transaminases (ALT/AST): Indicates that the transaminase elevation is not associated with significant hepatocyte necrosis/regeneration sufficient to produce an AFP elevation. Most viral hepatitis flares are associated with AFP <50--100 ng/mL that tracks with ALT
\end{itemize}

\subsection*{Above Normal}
\begin{itemize}
  \item AFP 20--200 ng/mL: Mild to moderate elevation. Common causes include benign liver diseases (acute hepatitis flare, chronic active HBV or HCV hepatitis, cirrhosis with active necroinflammation) where AFP tracks with ALT elevation. In the absence of active hepatitis, this range raises concern for early HCC
  \item AFP 200--400 ng/mL: Moderate elevation. In a patient with cirrhosis without an acute hepatitis flare, this level is highly suspicious for HCC and mandates liver imaging (multiphasic CT or MRI with LI-RADS assessment). Benign liver disease rarely produces sustained AFP >200 ng/mL without corresponding ALT >200--300 U/L
  \item AFP >400 ng/mL: High concern in a patient with cirrhosis or chronic hepatitis B, but AFP is not sufficient by itself to diagnose HCC. Use appropriate multiphasic imaging/LI-RADS or pathology; interpret AFP as surveillance, risk, or prognostic context.
  \item AFP >1,000 ng/mL: Strongly suggests HCC (in the context of cirrhosis or chronic HBV) or NSGCT (in the context of a testicular mass). IGCCCG criteria for NSGCT: AFP >1,000 ng/mL defines intermediate-risk disease; AFP >10,000 ng/mL defines poor-risk disease
  \item AFP >10,000 ng/mL: Almost always indicates either advanced HCC with a large tumor burden or aggressive NSGCT with yolk sac tumor components
  \item AFP elevation in NSGCT: AFP is exclusively produced by yolk sac (endodermal sinus) tumor elements and embryonal carcinoma with yolk sac differentiation. AFP is NEVER elevated in pure seminoma (or germinoma), pure choriocarcinoma (which produces only hCG), or pure mature teratoma. In a patient with a testicular mass and elevated AFP, there is unequivocally a non-seminomatous component; the tumor should be managed as NSGCT, not seminoma, even if histology appears seminomatous
  \item AFP-L3\% elevation: $\geq$10\% of total AFP in the L3 fraction is associated with HCC. AFP-L3\% has higher specificity (92--95\%) for HCC than total AFP alone. AFP-L3\% >10\% in a patient with total AFP 20--200 ng/mL enhances specificity and can detect HCC at an earlier stage. AFP-L3\% correlates with aggressive tumor biology (poor differentiation, portal vein invasion) and is a negative prognostic factor
  \item AFP elevation from non-malignant conditions: Acute and chronic hepatitis (AFP typically <200 ng/mL and correlates with ALT), cirrhosis (AFP 20--200 ng/mL), massive hepatic necrosis, drug-induced liver injury, pregnancy (physiologic elevation at 12--32 weeks), and hereditary persistence of AFP (a benign autosomal dominant condition)
  \item AFP elevation in non-HCC/GCT malignancies: Intrahepatic cholangiocarcinoma (occasional, typically <200 ng/mL), gastric adenocarcinoma with hepatoid differentiation, and hepatocellular adenoma (rare)
\end{itemize}

\section*{Normal Results}
\begin{itemize}
  \item \textbf{AFP (adults):} <10 ng/mL (or <8.5 ng/mL, laboratory-dependent)
  \item \textbf{AFP (borderline):} 10--20 ng/mL (equivocal; clinical significance uncertain in the absence of risk factors for HCC)
  \item \textbf{AFP (newborns):} AFP is markedly elevated at birth (30,000--120,000 ng/mL in term infants) and falls progressively, reaching adult levels (<10 ng/mL) by 10--12 months of age
  \item \textbf{Maternal serum AFP (pregnancy) --- median multiples of the median (MoM):}
  \begin{itemize}
    \item 15--20 weeks gestation: 0.5--2.5 MoM is normal
    \item $>$2.5 MoM: Elevated; associated with fetal neural tube defects (open spina bifida, anencephaly), abdominal wall defects (gastroschisis, omphalocele), multiple gestation, and fetal demise
    \item <0.5 MoM: Low; associated with trisomy 21 (Down syndrome) and trisomy 18 (Edwards syndrome) when combined with abnormal beta-hCG and inhibin A
  \end{itemize}
  \item \textbf{AFP-L3\%:} <10\% (normal). The AFP-L3\% assay is valid only when total AFP is $>$10 ng/mL
  \item \textbf{Hereditary persistence of AFP:} AFP 10--200 ng/mL in healthy individuals without liver disease or malignancy. Diagnosed by documenting elevated AFP in a first-degree relative
\end{itemize}

\section*{Follow-up Tests}
\begin{itemize}
  \item \textbf{Liver ultrasound:} The primary HCC screening modality. Should be performed every 6 months in patients with cirrhosis (regardless of etiology) per AASLD and EASL guidelines. Nodules <1 cm on ultrasound should be followed with repeat ultrasound in 3--6 months
  \item \textbf{Multiphasic contrast-enhanced CT of the liver (LI-RADS) or contrast-enhanced MRI of the liver:} For evaluation of a liver mass identified on ultrasound or suspected based on rising AFP. HCC characteristically shows arterial phase hyperenhancement followed by washout in the portal venous or delayed phase (LI-RADS LR-5 = definitely HCC, does not require biopsy for diagnosis)
  \item \textbf{AFP-L3\% (lectin-reactive AFP):} Measurement of the fucosylated isoform of AFP. Higher specificity for HCC than total AFP. Especially useful in patients with total AFP 10--200 ng/mL to distinguish benign liver disease from early HCC
  \item \textbf{PIVKA-II (protein induced by vitamin K absence II, also known as des-gamma-carboxy prothrombin or DCP):} An alternative tumor marker for HCC used in Japan and Asia. PIVKA-II is elevated in 50--70\% of HCC and is complementary to AFP --- combined AFP and PIVKA-II improves sensitivity for early HCC detection compared to either alone. Not widely used in North America or Europe
  \item \textbf{Liver biopsy (percutaneous, transjugular, or EUS-guided):} For histologic confirmation of HCC when imaging is not diagnostic (LI-RADS LR-3 or LR-4). Core needle biopsy preferred over fine-needle aspiration for adequate histologic assessment including tumor grading and immunohistochemistry
  \item \textbf{Beta-hCG (human chorionic gonadotropin) and LDH:} For staging, risk classification, and post-treatment monitoring of testicular germ cell tumors. IGCCCG classification: good-risk NSGCT (AFP <1,000 ng/mL, hCG <5,000 IU/L, LDH <1.5 $\times$ ULN), intermediate-risk, and poor-risk (any marker >10 $\times$ ULN, or primary mediastinal NSGCT)
  \item \textbf{Scrotal ultrasound and orchiectomy:} For diagnosis of testicular germ cell tumors. Radical inguinal orchiectomy provides the primary tumor for histologic classification. Trans-scrotal biopsy is contraindicated due to risk of tumor seeding
  \item \textbf{CT chest/abdomen/pelvis:} For staging of testicular germ cell tumors to detect retroperitoneal lymph node and pulmonary metastases
\end{itemize}

\section*{References}
\bibliographystyle{unsrtnat}
\bibliography{references}

\clearpage
\section*{Regulatory Status and Authorization Date}
This chapter describes a test category, not a specific commercial product. FDA, CDSCO, EU IVDR, and MHRA approval or registration dates therefore cannot be assigned without the assay manufacturer, product name, instrument, and intended use. For laboratory-developed tests, an individual agency approval date may not exist. See the Preface for the interpretation of regulatory status.

\clearpage
